\chapter{Software maintenance}

\section{Coding guidelines}

This program has been implemented following the Python PEP 8 style guide and targets Python 3.11
or newer; the exact supported range is declared in \keyw{pyproject.toml}.
The code has been documented using standard Python docstrings and type hinting.

Code quality is enforced automatically through a \keyw{pre-commit} configuration
(\keyw{.pre-commit-config.yaml}) that runs locally on the developer's machine before each commit.
The key checks include (see \keyw{.pre-commit-config.yaml} for the complete set):

\begin{itemize}
\item \emph{Black} for code formatting, with a line length of 88 characters that is shared with
\keyw{isort} and \keyw{flake8} (the \keyw{[tool.black]}, \keyw{[tool.isort]} and \keyw{[tool.flake8]}
sections of \keyw{pyproject.toml});
\item \emph{isort} for import ordering;
\item \emph{flake8} (with the \keyw{flake8-bugbear} and \keyw{flake8-docstrings} plugins) for linting
and PEP~257 docstring conventions; new docstrings follow the Google style;
\item \emph{bandit}, \emph{gitleaks}, \emph{detect-secrets}, \emph{trufflehog}, \emph{checkov} and
\emph{detect-private-key} for security and secret scanning;
\item \emph{beautysh} and \emph{shell-lint} for the shell and \keyw{.bat} build scripts;
\item \emph{commitizen} together with a set of commit-message hooks (imperative mood, capitalisation,
length and punctuation of the summary, blank second line) to enforce the Conventional Commits format;
\emph{commitizen} also drives version bumping and change-log generation (the \keyw{[tool.commitizen]}
section, \keyw{update\_changelog\_on\_bump});
\item a set of file-hygiene hooks (trailing whitespace, end-of-file, mixed line endings, large-file
guard, merge-conflict markers, JSON/YAML/TOML validation) and a \keyw{no-commit-to-branch} guard that
blocks direct commits to \keyw{main};
\item \emph{pytest} (unit, doctest and notebook runs) executed as local hooks, excluding the
end-to-end and binary tests.
\end{itemize}

The formatting and linting settings are defined in \keyw{pyproject.toml} (the \keyw{[tool.black]},
\keyw{[tool.isort]} and \keyw{[tool.flake8]} sections).
To run the full set of checks manually, install \keyw{pre-commit} and execute:

\begin{Verbatim}
    > pre-commit run --all-files
\end{Verbatim}

The \keyw{pre-commit} hooks are not run by continuous integration; they are intended to be installed
and run on the developer's machine. Continuous integration (the \keyw{testing.yml} GitHub Actions
workflow) instead runs the \keyw{pytest} test suite (excluding the binary tests) across a matrix of
Python 3.11--3.13 on Linux and Windows using Poetry, and uploads the coverage results to Codecov.
A minimum overall line-coverage threshold of 61\% is enforced; see \autoref{Chp:TestPlan} for details
of the test suite.

\section{Version control}

GitHub is currently used for software version control.
The repository is located at \url{https://github.com/Deltares/D-FAST_Bank_Erosion}.
Since \dfastbe builds on WAQBANK, the initial release of the new Python product is labeled as version 2.0.0.
%We may switch to GitLab in line with \href{http://publicaties.minienm.nl/documenten/rijkswaterstaat-informatievoorziening-aansluitvoorwaarden-riva-2017}{RIVA (2020)}.

\section{Automated building of code}

The standalone Windows executable is built locally with the \emph{Nuitka} compiler, driven by the
scripts in \keyw{BuildScripts} (\keyw{BuildDfastbe.bat}, or \keyw{BuildDfastbe\_no\_command\_window.bat}
for a build without a console window). The script reads the version number from \keyw{poetry version}
and invokes Nuitka roughly as follows:

\begin{Verbatim}
python -m nuitka --standalone --mingw64 --enable-plugin=pyside6 --enable-plugin=anti-bloat
    --python-flag=no_site --python-flag=no_asserts --python-flag=no_docstrings
    --include-package=numpy --include-package=pyproj --include-module=shapely
    --include-package=matplotlib --include-package=netCDF4 --include-package=cftime
    --include-module=geopandas --include-package=pandas --include-package=pytz
    --include-package-data=geopandas.datasets --include-module=fiona
    --include-data-files=src/dfastbe/io/log_data/messages.UK.ini=dfastbe/io/log_data/messages.UK.ini
    --file-reference-choice=runtime src/dfastbe
\end{Verbatim}

Unfortunately, this doesn't automatically build a binary that works out of the box.
The following adjustments depend on the exact combination of Nuitka and package versions used, so they will need to be revisited when updating the software.
A number of environmental settings (together with an explicit \keyw{pyproj} data-directory setup) had
to be added to \keyw{\_\_main\_\_} so that the standalone executable locates its bundled data
directories at runtime; these are settings that the Nuitka compiler did not resolve automatically.

\begin{Verbatim}
import os
from pathlib import Path

import matplotlib

matplotlib.use("Qt5Agg")

is_nuitka = "__compiled__" in globals()
if is_nuitka:
    """Needed for Nuitka compilation"""

    root = str(Path(__file__).parent)
    os.environ["GDAL_DATA"] = root + os.sep + "gdal"
    os.environ["PROJ_LIB"] = root + os.sep + "proj"
    os.environ["MATPLOTLIBDATA"] = root + os.sep + "matplotlib" + os.sep + "mpl-data"
    os.environ["TCL_LIBRARY"] = root + os.sep + "lib" + os.sep + "tcl8.6"
    proj_lib_dirs = os.environ.get("PROJ_LIB", "")
    import pyproj.datadir

    pyproj.datadir.set_data_dir(root + os.sep + "proj")
    import pyproj
\end{Verbatim}

If these environment settings are not applied, the standalone executable fails to locate its bundled
data directories at runtime, producing errors such as a missing PROJ, matplotlib or Tcl/Tk data
directory.

The Nuitka options above already bundle the required packages together with the data files that the
tool needs at runtime: the language files \keyw{messages.NL.ini} and \keyw{messages.UK.ini}, the GUI
icons, the licence file and the three PDF manuals.

A few shared libraries that Nuitka places in separate \keyw{*.libs} folders still have to be moved
next to the packages that load them. This is handled automatically by \keyw{BuildScripts/Move\_Libs.bat},
which moves the contents of \keyw{pyproj.libs}, \keyw{matplotlib.libs} and \keyw{fiona.libs} into the
corresponding \keyw{pyproj}, \keyw{matplotlib} and \keyw{fiona} folders of the distribution:

\begin{Verbatim}
move dfastbe.dist\pyproj.libs\*     dfastbe.dist\pyproj
move dfastbe.dist\matplotlib.libs\* dfastbe.dist\matplotlib
move dfastbe.dist\fiona.libs\*      dfastbe.dist\fiona
\end{Verbatim}

Finally, \keyw{BuildScripts/Collect\_Examples.bat} copies the example data sets into the distribution.
At runtime, the environment settings shown above point \keyw{pyproj}, GDAL, PROJ, matplotlib and
Tcl/Tk to the corresponding data directories inside the bundled distribution, so no further manual
copying of data directories or DLLs is required.

\section{Listing of external modules}

The runtime and development dependencies of \dfastbe --- including the supported Python version range
and the pinned version constraints --- are declared in \keyw{pyproject.toml}.
Refer to that file (and the accompanying lock file) for the authoritative, up-to-date list of
external modules rather than a static snapshot that quickly becomes outdated.

The main runtime dependencies are \keyw{numpy}, \keyw{pandas}, \keyw{geopandas}, \keyw{shapely},
\keyw{netCDF4}, \keyw{matplotlib}, \keyw{Pillow}, \keyw{pyside6} and \keyw{dfastio}.
Additional development, documentation and binary-build dependencies are grouped under the
\keyw{[dependency-groups]} section of \keyw{pyproject.toml}.

\section{Automated testing of code}

See \autoref{Chp:TestPlan} and \autoref{Chp:TestReport}.

\section{Automated generation of documentation}

The documentation has been written in a combination of LaTeX and markdown files which are maintained in the GitHub repository alongside the source code.
The PDF version of the user manual and this technical reference manual are generated automatically as part of the daily cycle of building all manuals on the Deltares TeamCity server.